\chapter{Całka Lebesgue'a w przestrzeni euklidesowej}
\label{ma2:ch:lebesgue-euclidean}

Celem rozdziału jest przygotowanie do obliczania
całek na zbiorach $\R^n$, a następnie na
krzywych, powierzchniach i podrozmaitościach.
W rozdziale \ref{ma2:ch:riemann-integral} zdefiniowano
całkę Riemanna na prostokątach i zawartość Jordana.
Całka Lebesgue'a umożliwia rozszerzenie tego
rachunku na dowolne zbiory mierzalne, funkcje
nieograniczone i granice ciągów funkcji.
Pełna konstrukcja miary i całki znajduje się w
\atlasref{ME-04::ch:lebesgue}{teorii miary i całki, rozdział 4}
oraz \atlasref{ME-07::ch:calkowanie-lebesgue}{rozdziale 7}.
Przytoczymy potrzebne definicje i twierdzenia,
nie powtarzając konstrukcji abstrakcyjnej.
W całej dalszej części rozpatrujemy funkcje skalarne;
pola wektorowe o skończenie wielu współrzędnych
pojawią się w zastosowaniach geometrycznych.

\section{Miara i całka na zbiorach w $\R^n$}

Niech $n\ge1$ i niech $\lambda_n$ oznacza
zupełną miarę Lebesgue'a na $\R^n$.
Dla prostokąta $Q=\prod_{j=1}^n[a_j,b_j]$,
gdzie $a_j\le b_j$, mamy
\[
 \lambda_n(Q)=\prod_{j=1}^n(b_j-a_j).
\]
Przyjmujemy standardowe własności miary
Lebesgue'a, w szczególności niezmienniczość
względem przesunięć, przeliczalną addytywność
i zupełność
(\atlasref{ME-04::thm:lebesgue-exists}{twierdzenie o istnieniu miary}).

\begin{definicja}
\label{ma2:def:lebesgue-domain-integral}
Niech $A\subseteq\R^n$ będzie mierzalny.
Dla mierzalnej funkcji nieujemnej
$f:A\to[0,\infty]$ definiujemy
\[
 \int_A f(x)\,\mathrm d\lambda_n(x)
 =\sup\left\{
 \sum_{j=1}^m c_j\lambda_n(A_j):
 \ \sum_{j=1}^m c_j\chi_{A_j}\le f,\quad
 c_j>0
 \right\},
\]
gdzie $A_j$ są parami rozłącznymi
mierzalnymi podzbiorami $A$, $m$ jest
skończone, a do supremum dołączamy wartość
$0$ odpowiadającą funkcji zerowej.
Dla mierzalnej funkcji rzeczywistej $f$
kładziemy $f^+=\max(f,0)$ i $f^-=\max(-f,0)$.
Jeśli $\int_A |f|\,\mathrm d\lambda_n<\infty$,
nazywamy $f$ całkowalną i definiujemy
\[
 \int_A f\,\mathrm d\lambda_n
 =\int_A f^+\,\mathrm d\lambda_n
  -\int_A f^-\,\mathrm d\lambda_n.
\]
\end{definicja}

W szczególności dla $B\subseteq A$ mierzalnego
\[
 \int_A\chi_B(x)\,\mathrm d\lambda_n(x)=\lambda_n(B).
\]
Wynika to z definicji, ponieważ sama
$\chi_B$ jest funkcją prostą, a całka funkcji
prostej jest monotoniczna
(\atlasref{ME-07::thm:simple-properties}{własności całek prostych}).
Całka funkcji nieujemnej może mieć wartość
$+\infty$; dla funkcji ze znakiem nie
dopuszczamy różnicy $\infty-\infty$.

\begin{stwierdzenie}[Przedłużenie przez zero]
\label{ma2:prop:lebesgue-zero-extension}
Dla mierzalnego $A\subseteq\R^n$ i mierzalnej
funkcji $f:A\to\R$ niech
\[
 \widetilde f(x)=
 \begin{cases}
 f(x),&x\in A,\\
 0,&x\notin A.
 \end{cases}
\]
Wówczas $\widetilde f$ jest mierzalna,
a całka $f$ na $A$ jest równa całce
$\widetilde f$ na $\R^n$, jeżeli $f$
jest nieujemna lub całkowalna.
\end{stwierdzenie}
\begin{proof}
Dla borelowskiego $B\subseteq\R$ przeciwobraz
$\widetilde f^{-1}(B)$ jest równy
$f^{-1}(B)$ lub $f^{-1}(B)\cup(\R^n\setminus A)$,
w zależności od tego, czy $0\in B$.
Oba zbiory są mierzalne, gdyż $A$
jest mierzalny, a mierzalność $f$
rozumiemy względem ograniczenia sigma-ciała do $A$.
Dla $f\ge0$ funkcje proste pod $\widetilde f$
są zerowe poza $A$, więc ich ograniczenia
do $A$ dają dokładnie te same wartości całek
prostych co w definicji
\ref{ma2:def:lebesgue-domain-integral}.
Odwrotnie, każdą funkcję prostą pod $f$
przedłużamy przez zero. Supremum jest więc
takie samo. Dla funkcji ze znakiem stosujemy
wynik do $f^+$ i $f^-$ osobno; równość dla
$|f|$ zapewnia również całkowalność przedłużenia.
\end{proof}

\section{Zgodność z całką Riemanna na zbiorach Jordana}

Niech $A$ będzie ograniczonym zbiorem
mierzalnym w sensie Jordana, a $Q$
prostokątem zawierającym $A$.
Dla ograniczonej funkcji $f:A\to\R$
całkę Riemanna na $A$ definiujemy jako
całkę przedłużenia przez zero na $Q$,
o ile ta istnieje. Niezależność od $Q$
wynika z addytywności całki Riemanna
na prostokątach: dla dwóch prostokątów
zawierających $A$ bierzemy większy
prostokąt zawierający oba i rozdrabniamy
go wzdłuż ścian każdego z nich.
Na pozostałych komórkach przedłużenie
jest zerowe poza ścianami, które mają
miarę zero, zatem całka na każdej z nich
jest zerowa z kryterium Lebesgue'a
i zgodności całek z rozdziału 24.

\begin{twierdzenie}
\label{ma2:thm:lebesgue-jordan-agreement}
Jeżeli $A\subseteq\R^n$ jest ograniczonym
zbiorem mierzalnym w sensie Jordana,
a $f:A\to\R$ jest ograniczona i ciągła
względnie na $A$, to $f$ jest całkowalna
w sensie Riemanna i Lebesgue'a na $A$,
a obie całki są równe.
\end{twierdzenie}
\begin{proof}
Wybieramy $Q$ zawierający domknięcie $A$
we wnętrzu. Przedłużenie przez zero
jest ciągłe w każdym punkcie wnętrza $A$,
gdyż w otoczeniu tego punktu pokrywa się
z $f$. Jest również ciągłe poza domknięciem
$A$, gdyż lokalnie jest zerowe.
Zatem zbiór jego punktów nieciągłości
jest zawarty w brzegu $\partial A$.
Z kryterium Jordana
\ref{ma2:thm:jordan-criterion}
brzeg ma miarę Lebesgue'a zero.
Przedłużenie jest ograniczone, więc
kryterium Lebesgue'a dla całki Riemanna
\ref{ma2:thm:riemann-box-lebesgue-criterion}
zapewnia jej istnienie. Z twierdzenia
\ref{ma2:thm:riemann-box-lebesgue}
przedłużenie jest mierzalne i jego całki
Riemanna i Lebesgue'a są równe.
Ograniczenie do mierzalnego $A$ jest
mierzalne, a stwierdzenie
\ref{ma2:prop:lebesgue-zero-extension}
daje żądaną zgodność na $A$.
\end{proof}

Wniosek ten pozwala obliczać całki
ciągłych funkcji na ograniczonych
zbiorach Jordana za pomocą dowolnego
z dwóch rachunków. W kolejnych
sekcjach rozważymy przechodzenie
do granicy i całki na zbiorach
niemierzalnych w sensie Jordana,
lecz mierzalnych w sensie Lebesgue'a.
